% ==============================================================================
%  Conviction-based inference — main subsection
%  Numbers come from quick_analysis/Conviction_experiments_tables/_summary/,
%  produced by quick_analysis/conviction_summary_tables.py with
%      --score last --exclude cornell,wisconsin
%  Full tables, robustness checks and procedural detail: conviction_supplementary.tex
% ==============================================================================

\subsection{Conviction-based inference}
\label{sec:conviction}

The Forman profile of Section~\ref{sec:forman_profile} is an observable of the
\emph{operator}: it describes the sheaf Laplacian that the network learns, and
says nothing about the signal that the operator propagates. The Bochner
decomposition, however, also gives us a reading of the dynamics itself. As we
discussed in Section~\ref{sec: boch}, sheaf diffusion can be interpreted as an
opinion dynamics model, in which each node carries an opinion and the
restriction maps encode how opinions are translated between neighbouring points
of view. In this reading the magnitude of a node's representation is not an
incidental quantity: it measures how strongly that node holds its position, and
we therefore call it the node's \emph{conviction}.

What makes this quantity interesting is that the theory tells us where it should
end up. For a Bochner Laplacian, that is for an unnormalised bundle sheaf
without edge weights, the equilibrium configuration requires all the stalk norms
to coincide (METTERE REF, TEOREMA). The conviction is therefore not merely a
descriptive statistic, but a quantity for which we have a theoretical
prediction, and whose deviation from that prediction is itself informative.

We use this in two ways. First \emph{observationally}, asking whether the
conviction that the network actually produces behaves as the equilibrium theory
predicts, and how it relates to simpler node properties such as degree. Second,
and this is the main contribution of this section, \emph{operationally}: we ask
whether conviction identifies the nodes on which the trained network relies, by
perturbing nodes of different conviction and measuring the damage. In this way a
quantity derived from the Bochner theory becomes an experimentally verifiable
score of node importance.

\subsubsection{Methods}

\paragraph{The conviction.}
In the sheaf setting the hidden state of node $u$ at diffusion layer $t$ is not
a vector but a matrix, $X_t[u] \in \mathbb{R}^{d \times F}$, where $d$ is the
stalk dimension and $F$ the number of hidden channels.

\begin{definition}[Conviction]
\label{def:conviction}
The \emph{conviction} of node $u$ at layer $t$ is the $L_2$ norm of its stalk,
\[
    c_t(u) \;=\; \bigl\| X_t[u] \bigr\|_{L_2}
    \;=\; \sqrt{\sum_{i=1}^{d}\sum_{j=1}^{F} \bigl(X_t[u]\bigr)_{ij}^2}
    \;\in\; \mathbb{R}_{\geq 0},
\]
and the \emph{conviction vector} is
$\mathbf{c}_t = (c_t(1),\dots,c_t(n)) \in \mathbb{R}^n$.
\end{definition}

We outline immediately a discrepancy with the theory that should be kept in
mind throughout. In the equilibrium result the state of a node is an element of
$\mathbb{R}^d$, whereas here it lives in $\mathbb{R}^{d \times F}$, since the
network also carries feature channels and interleaves the diffusion with weight
matrices. The setting is therefore richer than the one in which the theorem is
proved, and the equilibrium prediction should be read as a guide rather than as
a statement that applies verbatim. We nonetheless find it meaningful, because
the quantity it constrains, the magnitude of a node's representation, is exactly
the one we measure.

Unless stated otherwise we use the conviction at the \textbf{last diffusion
layer}, $\mathbf{c}_L$, taken at the best-validation-epoch checkpoint and
denoted \texttt{last}. As a robustness check we also considered
\texttt{avg\_z}, obtained by z-scoring the conviction of each post-diffusion
layer $1,\dots,L$ separately and averaging, so that no single layer's scale can
dominate; the conclusions below are unchanged under this choice, and the
corresponding tables are reported in the supplementary material.

\paragraph{Models and datasets.}
We study two model families. \textbf{GeneralSheaf} with
$\texttt{normalised=True}$ learns unconstrained $d \times d$ restriction maps
and is the configuration of interest from the machine learning standpoint, being
the most expressive of the sheaf learners we consider. \textbf{BundleSheaf} with
$\texttt{normalised=False}$ and no edge weights learns orthogonal restriction
maps through a Householder parametrisation, and is the configuration of interest
from the mathematical standpoint, since its Laplacian is a Bochner Laplacian and
the equilibrium prediction applies to it.

The experimental protocol matches Section~\ref{sec:forman_profile}: the same ten
node classification benchmarks, $10$-fold cross-validation, $500$ epochs with
early stopping at patience $200$, and per-dataset choices of stalk dimension
$d \in \{2,3,5\}$ and depth $L \in \{2,\dots,6\}$. Performance is test accuracy,
and ROC-AUC for Minesweeper, Tolokers and Questions, which are binary and
strongly imbalanced.

\paragraph{Observational analysis.}
For each dataset and model we compare the conviction against other node
properties at every training checkpoint and every diffusion layer, producing a
grid of heatmaps whose rows are checkpoints and whose columns are layers. Each
cell reports one scalar averaged over folds: the Spearman correlation of
$\mathbf{c}_t$ with node degree, with betweenness centrality, with the model's
softmax confidence and with input-gradient saliency, together with the absolute
Cohen's $d$ separating the conviction of correctly and incorrectly classified
nodes, computed separately on train and test. We use the absolute value of
Cohen's $d$ before averaging over folds, so that sign cancellations across folds
do not hide an effect that is present in each of them. In addition, at the
best-epoch checkpoint we plot conviction against degree, one panel per diffusion
layer, to inspect directly whether high-conviction nodes are simply hubs.

\paragraph{From observation to intervention.}
The observational analysis can only tell us what conviction correlates with. To
establish that it identifies nodes the network actually depends on, we move to
an interventional setting. The protocol is deliberately \emph{inference-only}:
we load the frozen best-epoch weights, perturb a selected set of nodes, and
re-evaluate. Nothing is retrained, so any change in performance is attributable
to the perturbation rather than to a different optimisation trajectory.

We use three perturbation modes, which act on different parts of the mechanism:
\begin{enumerate}
    \item \textbf{masking}, which sets the selected nodes' input features to
    zero, or replaces them with the dataset feature mean; the topology is
    untouched, so this probes the reliance on those nodes' \emph{features};
    \item \textbf{gating}, which multiplies the selected nodes' outgoing
    messages by zero while leaving their self-term and every other node's
    messages intact, so this probes the reliance on what those nodes
    \emph{send};
    \item \textbf{pruning}, which removes the selected nodes together with their
    incident edges and forward-passes on the smaller graph, reloading the same
    frozen weights into a model shaped for the new graph.
\end{enumerate}
Using three modes is not redundancy: an effect that survives all three cannot be
an artefact of one particular way of perturbing the network.

\paragraph{Selection strategies.}
Candidates are always drawn from the \textbf{training nodes only}, so that the
test set on which we evaluate is identical across every strategy and every
budget, and no test node is ever perturbed. At each budget
$k \in \{1\%, 5\%, 10\%, 20\%, 30\%\}$ of the training pool we compare nine
selection strategies: the top and bottom $k\%$ by conviction, by degree and by
prediction confidence; $k$ uniformly random nodes; and two
\emph{degree-matched} controls, in which $k$ nodes are sampled to reproduce the
degree distribution of the high- or low-conviction group while being otherwise
random.

The degree-matched controls are the decisive comparison of this section. Degree
is the obvious competing explanation for any effect we might find, and matching
the degree distribution while randomising the identity of the nodes isolates
exactly the information that conviction carries \emph{beyond} topology. Ties in
any ranking, which are extremely common for degree, are broken by a per-node
random key seeded deterministically from the fold, so that the selection does
not systematically favour low-index nodes and is invariant to the arbitrary
labelling of the graph.

\paragraph{Decile analysis.}
The budget sweep is cumulative: at $k\%$ we perturb the top $k\%$ all at once,
so the high-conviction nodes of the $1\%$ budget are contained in every larger
one. This tells us that high conviction matters, but not \emph{how} the damage
is distributed across the conviction range. We therefore also run a decile
experiment, in which the training pool is sorted by a criterion, split into ten
equal bins, and \textbf{exactly one bin is perturbed at a time}. Every bin has
the same size, so the ten resulting numbers are directly comparable, and they
answer a sharper question: is conviction a binary high/low label, or a
calibrated ordering along which importance grows smoothly? We run the same
decile experiment for conviction, for confidence and for degree, since the
comparison between the three is the entire point.

\paragraph{Aggregation.}
Reporting ten datasets $\times$ nine strategies $\times$ five budgets $\times$
three perturbation modes separately would be unreadable. We therefore express
every result as a \emph{relative drop} with respect to that dataset's own
unperturbed baseline,
\[
    \mathrm{drop}\%(s) \;=\; 100 \cdot
    \frac{\overline{m}_{\text{baseline}} - \overline{m}_{s}}
         {\overline{m}_{\text{baseline}}},
\]
where $\overline{m}$ denotes the headline metric already averaged over the ten
folds, and we then average over datasets. Each dataset contributes one value and
counts equally, so that the large graphs do not dominate the aggregate, and the
quoted uncertainty is the spread \emph{across datasets}, which is what tells us
whether an effect would reproduce on a new benchmark.

We exclude Cornell and Wisconsin from the aggregation. Their training pools
contain fewer than $130$ nodes, so a $1\%$ budget is a single node and a decile
fewer than ten, and the resulting values are dominated by sampling noise while
entering the average with the same weight as a graph with tens of thousands of
nodes. Including them does not change the ordering of the strategies, only its
sharpness; we report the eight-dataset aggregate throughout, and give the
ten-dataset version in the supplementary material.

Because the relative drops are strongly right-skewed across datasets, we do not
rely on the mean alone. We report, per strategy, the mean rank within each
dataset (rank $1$ = most damaging), and we test the comparisons the section
turns on with a Wilcoxon signed-rank test in which \textbf{the pairing unit is
the dataset}, following the standard practice for comparing methods over a
benchmark suite (METTERE REF: Dem\v{s}ar, 2006). We stress that the test is
across datasets and not across folds: the per-dataset tables store fold
averages, and the baseline and perturbed metrics share the same folds, so
propagating fold-level errors into the ratio would ignore that correlation.

\subsubsection{Results}

\paragraph{The bundle sheaf does not reach the predicted equilibrium.}
We begin with the observational analysis of the unnormalised bundle sheaf, the
configuration for which the theory makes a prediction. If the learnt dynamics
converged to the equilibrium of the Bochner Laplacian, the conviction at the
final diffusion layer would be constant across nodes. This is emphatically not
what we observe: the conviction at the last layers remains spread over a wide
range of values, and the nodes retain markedly different magnitudes. In
addition, the conviction of the bundle sheaf is strongly correlated with node
degree and with betweenness centrality, so the surviving structure in the norms
is largely topological.

We regard this as the mathematically interesting finding of the observational
part. The system is \textbf{not at equilibrium}, and consequently the standard
kernel-theoretic machinery used to study the convergence and the limit points of
sheaf diffusion, or of graph convolutional networks, does not describe what the
trained network is doing. The reason, in our reading, is the same one that
motivated the Forman profile: a Sheaf Neural Network is not only a diffusion,
since the sheaf itself is re-estimated at every layer from the current features.
It is a \textbf{geometric flow}, and a geometric flow has no reason to settle at
the equilibrium of the frozen diffusion associated with any one of its
instantaneous operators. The narrative according to which representations
converge to a common value is therefore perturbed, and the observed
out-of-equilibrium behaviour is a direct empirical signature of the flow.

We repeat the caveat of the Methods here, because it bears on how strongly this
should be read: our states live in $\mathbb{R}^{d\times F}$ rather than
$\mathbb{R}^d$, and the network interleaves weight matrices with the diffusion,
so the equilibrium theorem does not apply verbatim. What we can say is that the
quantity the theorem constrains does not behave as the theorem would suggest,
and that the geometric-flow reading accounts for this naturally.

\paragraph{In the general sheaf, conviction is not degree.}
For the normalised general sheaf the picture is different, and it is this
configuration that matters for the machine learning question. The conviction is
only weakly correlated with degree and with betweenness, and inspecting the
scatter plots layer by layer confirms that the ordering induced by conviction
does not coincide with the ordering induced by degree; Roman Empire, Questions
and Pubmed are clear examples, and the same holds for the majority of the
datasets we consider. Conviction is instead appreciably correlated with the
model's softmax confidence, while its correlation with input-gradient saliency
is weak and sometimes negative. The Cohen's $d$ separating correctly from
incorrectly classified nodes is appreciable at early epochs and becomes small
later in training.

Since conviction is neither a restatement of degree nor a restatement of
confidence, the question of whether it identifies something the network actually
uses is a genuine one, and it cannot be settled by correlations. This is what
motivates the interventional experiments.

\paragraph{High-conviction nodes are the ones the trained network relies on.}
The central result is summarised in Table~\ref{tab:conviction_forward} and is
remarkably stable. Across all three perturbation modes and at every budget,
removing the highest-conviction nodes is the most damaging thing one can do, and
removing the lowest-conviction nodes is the least damaging. Under pruning at a
$30\%$ budget, for instance, the high-conviction selection costs $4.93\%$ of the
metric while the low-conviction selection costs $0.55\%$, an order of magnitude
less, with random removal at $1.97\%$ sitting between them.

The margin over the controls is statistically reliable. Against random
selection, against the degree-matched control, against confidence-based
selection and against low-conviction selection, the high-conviction strategy is
more damaging in seven or eight of the eight datasets at essentially every
budget and in every perturbation mode, with Wilcoxon $p \leq 0.016$ throughout.
The degree-matched comparison deserves emphasis: a set of nodes reproducing the
degree distribution of the high-conviction group, but otherwise random, costs
only $0.65\%$ under pruning at $10\%$, against $2.70\%$ for the high-conviction
group itself. Conviction is therefore \textbf{not} a proxy for degree; knowing
how connected the nodes are is not enough, one has to know which nodes they are.

We must however be precise about one comparison, since it constrains the claim.
Selecting the top-$k\%$ nodes \emph{by degree} is also a strong strategy, and
although high conviction has the larger mean drop in every mode, the difference
between the two is not always reliable. It becomes significant from the $5\%$
budget onwards for hard gating and for zero-masking, and at the largest budgets
for mean-masking, but it is never significant under pruning. This last exception
has a natural mechanistic explanation: pruning deletes nodes together with their
incident edges, so it is precisely the mode in which a degree-based criterion is
maximally advantaged, since degree is by definition the number of edges removed.
We therefore do not claim that conviction outperforms degree as a selection
rule at a fixed budget; we claim that it is the most damaging criterion overall,
that it is reliably better than every control that is not degree itself, and
that it carries information beyond the degree distribution.

The mirror-image statement is equally useful and is the one we would emphasise
for applications. Low-conviction removal is significantly \emph{less} damaging
than random removal at the larger budgets, in seven or eight of eight datasets
($p \leq 0.039$ under pruning from the $5\%$ budget onwards). Low-conviction
nodes are thus not merely unimportant, they are demonstrably more dispensable
than an arbitrary node would be.

\begin{table}[H]
\centering
\small
\caption{Relative drop (\%) in the headline metric under node pruning, as a
function of the fraction of the training pool perturbed. Conviction score
\texttt{last}; mean over the eight datasets retained, with the standard error of
the cross-dataset mean in parentheses. Larger is more damaging. The equivalent
tables for the other two perturbation modes are given in the supplementary
material.}
\label{tab:conviction_forward}
\begin{tabular}{lccccc}
\toprule
Selection & 1\% & 5\% & 10\% & 20\% & 30\% \\
\midrule
High conviction        & $0.91$ & $1.88$ & $2.70$ & $3.92$ & $4.93$ \\
High degree            & $0.79$ & $1.68$ & $2.28$ & $3.32$ & $4.02$ \\
High confidence        & $0.24$ & $0.73$ & $1.12$ & $2.02$ & $2.81$ \\
Degree-matched (high)  & $0.16$ & $0.40$ & $0.65$ & $0.96$ & $1.27$ \\
Random                 & $0.08$ & $0.30$ & $0.60$ & $1.20$ & $1.97$ \\
Low confidence         & $0.07$ & $0.28$ & $0.54$ & $1.21$ & $1.90$ \\
Degree-matched (low)   & $0.04$ & $0.27$ & $0.45$ & $1.23$ & $2.12$ \\
Low degree             & $0.03$ & $0.16$ & $0.38$ & $0.78$ & $1.32$ \\
Low conviction         & $0.01$ & $0.06$ & $0.12$ & $0.29$ & $0.55$ \\
\bottomrule
\end{tabular}
\end{table}

\paragraph{Conviction grades importance, it does not merely label it.}
The decile experiment sharpens the previous paragraph considerably, and is
reported in Table~\ref{tab:conviction_decile}. Perturbing a single decile at a
time, the damage caused by the conviction bins grows monotonically from the
bottom of the distribution to the top, spanning more than an order of magnitude:
under pruning it rises from $0.12\%$ for the lowest decile to $2.70\%$ for the
highest. The comparison with the other two criteria is what makes this
informative. Confidence produces an almost flat profile, from $0.54\%$ to
$1.12\%$, a factor of two across the entire range; degree grows, but from a
higher floor and to a lower ceiling.

Conviction therefore has simultaneously the \emph{least} damaging bottom bin and
the \emph{most} damaging top bin of the three criteria, which is precisely what
it means for a score to be well calibrated. It is not that conviction separates
an important group from an unimportant one, a binary distinction that degree
also achieves; it is that the position of a node along the conviction range
predicts how much the network depends on it. This is the sense in which the
Bochner theory yields an operational quantity, and it is also the place where
conviction distinguishes itself most clearly from the degree baseline that
competes with it in the fixed-budget experiment.

\begin{table}[H]
\centering
\small
\caption{Decile analysis under node pruning: relative drop (\%) in the headline
metric when a single bin of the training pool is removed, the pool being split
into ten equal bins ranked by each criterion (bin $1$ = lowest score, bin $10$ =
highest). Conviction score \texttt{last}; mean over the eight datasets retained.
Only conviction produces a profile that grows by more than an order of magnitude
from the bottom bin to the top one.}
\label{tab:conviction_decile}
\begin{tabular}{lccc}
\toprule
Bin & Conviction & Confidence & Degree \\
\midrule
1  & $0.12$ & $0.54$ & $0.37$ \\
2  & $0.15$ & $0.61$ & $0.40$ \\
3  & $0.20$ & $0.57$ & $0.44$ \\
4  & $0.36$ & $0.57$ & $0.40$ \\
5  & $0.51$ & $0.66$ & $0.49$ \\
6  & $0.58$ & $0.61$ & $0.48$ \\
7  & $0.63$ & $0.65$ & $0.54$ \\
8  & $0.75$ & $0.62$ & $0.54$ \\
9  & $0.94$ & $0.69$ & $0.79$ \\
10 & $2.70$ & $1.12$ & $2.30$ \\
\bottomrule
\end{tabular}
\end{table}

\paragraph{Soft gating is inconclusive.}
For completeness we also ran a \emph{soft} gating variant, in which no subset is
selected and instead every training node's outgoing message is attenuated
continuously in proportion to its own score, the lowest-scoring node being
suppressed the most. This yields a single number per fold rather than a budget
sweep. The result is not interpretable in the same way as the others:
attenuating by conviction costs $2.44 \pm 0.96\%$, by degree $3.72 \pm 1.99\%$
and by confidence $0.50 \pm 0.36\%$, with error bars so wide that the three
cannot be separated. This is expected in hindsight, since the soft variant
simultaneously suppresses the low-conviction nodes and preserves the
high-conviction ones, so the two effects that the budget sweep isolates are here
superimposed and partly cancel. We report it as a negative result and base no
claim on it.

\subsubsection{Discussion}

The experiments of this section change the status of the conviction. It enters
as a descriptive quantity, read off the Bochner decomposition and constrained by
an equilibrium theorem; it leaves as an \textbf{operational score of node
importance}, whose predictions about which nodes matter are verified by
intervention on a trained network.

We would state the conclusion conservatively, because the conservative version
is the one the data supports and it is already a strong claim. High-conviction
nodes are those on which the learnt sheaf diffusion relies most, and
low-conviction nodes are correspondingly \emph{dispensable} for the already
trained model. What we do \emph{not} claim is that removing low-conviction nodes
improves the model. An inspection of the individual masking curves suggested a
slight improvement when low-conviction nodes were suppressed, but once the
results are aggregated across datasets with error bars this effect does not
survive: the honest statement is that low-conviction removal is the least
destructive intervention available, not that it is beneficial. This is a more
defensible claim, and it is also the more useful one, since it is what would
justify using conviction for interpretability, for robustness analysis, or for
compressing a trained model.

Two properties make the result more than an accuracy comparison. The first is
its robustness across perturbation modes: masking acts on features and leaves
the graph intact, gating acts on messages and leaves both features and graph
intact, pruning acts on the topology, and the ordering of the strategies is the
same in all three. An artefact of one mechanism would not survive this. The
second is the degree-matched control, which rules out the most natural
deflationary explanation. Conviction correlates only weakly with degree in the
general sheaf, and nodes reproducing the degree distribution of the
high-conviction set are far less important than the high-conviction set itself.

The honest limitation is the comparison with degree as a selection rule, which
we have stated in the Results and do not wish to bury. At a fixed budget,
sorting by degree is competitive with sorting by conviction, and under pruning
the two cannot be distinguished at all. What separates them is not the
fixed-budget experiment but the decile one, where conviction alone produces a
graded profile spanning an order of magnitude. The value of conviction, on this
evidence, is that it is \emph{calibrated}, not that it is the sharpest possible
way of finding the single most important node.

Finally, the observational half of this section connects back to
Section~\ref{sec:forman_profile}, and the two support the same reading from
opposite directions. There we found that the learnt Laplacian is not the object
the static theory would predict, and here we find that the signal it propagates
does not reach the equilibrium the static theory would predict either. Both are
consequences of the same fact: the sheaf is re-estimated at every layer, and the
network therefore deforms its own geometry as it diffuses. That the resulting
out-of-equilibrium conviction turns out to be a useful importance score is, we
think, the most interesting aspect of the result, since it suggests that what
the static theory treats as a failure to converge is in fact where the useful
structure resides.
